
% 1. Encoding & Fonts (基础设置)
\usepackage[utf8]{inputenc}
\usepackage[T1]{fontenc}
\usepackage{microtype}

% 2. Colors (必须合并加载，否则会报 Option clash)
% 将原本分散的三行 xcolor 合并为一行，一次性传入所有参数
\usepackage[table,dvipsnames]{xcolor}

% 3. Math Packages (数学工具)
\usepackage{amsmath}
\usepackage{amssymb}
\usepackage{mathtools}
\usepackage{amsthm}
\usepackage{amsfonts}
\usepackage{nicefrac}
\usepackage{bm}

% 4. Tables (表格工具)
\usepackage{booktabs}       % 专业表格线
\usepackage{multirow}
\usepackage{threeparttable}
\usepackage{colortbl}
\usepackage{array}
\usepackage{tabularx}

% 5. Figures & Graphics (图形工具)
\usepackage{graphicx}
\usepackage{wrapfig}
\usepackage{stfloats}
% === 关键修改：只保留 subcaption，删除了 subfigure ===
\usepackage{subcaption} 
\captionsetup[subfigure]{skip=0.5em}

% 6. Algorithms (ICML 标准方案)
% -----------------------------------------------------------
\makeatletter
% algorithm / endalgorithm
\@ifundefined{algorithm}{}{\let\algorithm\relax}
\@ifundefined{endalgorithm}{}{\let\endalgorithm\relax}

% algorithm* / endalgorithm*  (关键：你现在报错的就是这个)
\expandafter\ifx\csname algorithm*\endcsname\relax\else
  \expandafter\let\csname algorithm*\endcsname\relax
\fi
\expandafter\ifx\csname endalgorithm*\endcsname\relax\else
  \expandafter\let\csname endalgorithm*\endcsname\relax
\fi

% listofalgorithms (你截图的第二个报错)
\@ifundefined{listofalgorithms}{}{\let\listofalgorithms\relax}

% algorithmic (有时也会冲突)
\@ifundefined{algorithmic}{}{\let\algorithmic\relax}
\makeatother
% -----------------------------------------------------------

\usepackage[linesnumbered,ruled,vlined]{algorithm2e}
\SetKwComment{Comment}{$\triangleright$\ }{}
\SetKwRepeat{Do}{do}{while}

% 如果你需要 Input/Output，用 algorithm2e 的方式：
\SetKwInOut{Input}{Input}
\SetKwInOut{Output}{Output}
% \SetNlSty{textbf}{}{:}
\SetNlSkip{0.5em}
\setlength{\algomargin}{1.2em}
\newcommand{\cmt}[1]{\tcp*[r]{#1}} % (A) 行尾右对齐注释：\cmt{...}
\newcommand{\step}[1]{\tcp{\textbf{#1}}} % (B) 分节/步骤注释（整行强调）：\step{...}

% 7. Misc
\usepackage{url}
\usepackage{tcolorbox}
\usepackage{enumitem}
\usepackage{balance}
\usepackage{textcomp}
\usepackage{verbatim}

% 8. Hyperref
\definecolor{ICMLBlue}{RGB}{25, 88, 166}
\definecolor{ICMLPurple}{RGB}{150, 0, 80}
\usepackage[colorlinks,linkcolor=ICMLPurple,citecolor=ICMLBlue,urlcolor=black]{hyperref}
% \usepackage{hyperref}

% 9. Cleveref
\usepackage[capitalize,noabbrev]{cleveref}

%%%%%%%%%%%%%%%%%%%%%%%%%%%%%%%%
% THEOREMS
%%%%%%%%%%%%%%%%%%%%%%%%%%%%%%%%
\theoremstyle{plain}
\newtheorem{theorem}{Theorem}[section]
\newtheorem{proposition}[theorem]{Proposition}
\newtheorem{lemma}[theorem]{Lemma}
\newtheorem{corollary}[theorem]{Corollary}
\theoremstyle{definition}
\newtheorem{definition}[theorem]{Definition}
\newtheorem{assumption}[theorem]{Assumption}
\theoremstyle{remark}
\newtheorem{remark}[theorem]{Remark}

% Todonotes is useful during development; simply uncomment the next line
%    and comment out the line below the next line to turn off comments
%\usepackage[disable,textsize=tiny]{todonotes}
\usepackage[textsize=tiny]{todonotes}
\newcommand{\mytodo}[1]{\textcolor{red}{#1}}
\renewcommand{\paragraph}[1]{\noindent\textbf{#1}}
\renewcommand{\subparagraph}[1]{\noindent\textbf{\underline{#1}}}

\newcommand{\Wm}{\bm{W}_m}
\newcommand{\Wmi}{\bm{W}_{m,i}}
\newcommand{\Wq}{\bm{W}_q}
% \newcommand{\Wk}{\bm{W}_{\text{k}}}
\newcommand{\Wk}{\bm{W}_\mathrm{k}}
\newcommand{\Wv}{\bm{W}_\mathrm{v}}
\newcommand{\Wo}{\bm{W}_\mathrm{o}}
\newcommand{\Wup}{\bm{W}_\mathrm{up}}
\newcommand{\Wgate}{\bm{W}_\mathrm{gate}}
\newcommand{\Wdown}{\bm{W}_\mathrm{down}}
\newcommand{\Wemb}{\bm{W}_\mathrm{emb}}
\newcommand{\Wlm}{\bm{W}_\mathrm{lm}}
\newcommand{\MSE}{\mathcal{E}}
\newcommand{\mb}[1]{\bm{#1}}
\newcommand{\mrm}[1]{\mathrm{#1}}
\newcommand{\mbb}[1]{\mathbb{#1}}
\newcommand{\mcal}[1]{\mathcal{#1}}
